% Sample Applications, Chapter 5 -- Monthly Expense Report (annex Sec 8.3).
% Values replayed through the shipped application (examples/expenses5.i23);
% @SUM and @AVG as os/i123/core/eval.c computes them, the Currency display
% as os/i123/core/fmt.c.  Figures screens/sa5.k123, sa5b.k123.
% GATES: test-i123-formula test-pubs-i123-examples test-pubs-screens
\Gates{test-i123-formula, test-pubs-i123-examples, test-pubs-screens}
\chapter{Monthly Expense Report}

At the end of each month the office manager lists what the office spent,
and reports the total and the average expense.  This worksheet adds up a
list of expenses and shows every amount in dollars and cents.

\section{What You Enter}
\begin{itemize}
\item A short name for each expense
\item The amount spent on it
\end{itemize}

\section{Entering the Data}
\ix{expense report}\ix{"@AVG}
Type the headings \type{Item} and \type{"Cost} in A1 and B1.  Type the
expenses \type{Rent}, \type{Phone}, \type{Travel} and \type{Fuel} in A2..A5
and the labels \type{Total} and \type{Average} in A6 and A7, then the
amounts in B2..B5.  In B6 enter \type{@sum(b2..b5)} and in B7 enter
\type{@avg(b2..b5)}.  Finally, format B2..B7 as currency with two decimal
places: select \slashcmd{/Range Format Currency}, type \type{2}, and give
the range.
\FlairShot[fig:sa5]{sa5}{The month's expenses with total and average}

\begin{LotTable}{p{30mm}p{24mm}p{70mm}}
\LotHead{\textbf{Item} & \textbf{Address} & \textbf{Meaning}}
Expense & A2..A5 & The name of each expense\\
Cost & B2..B5 & What was spent on it\\
Total & B6 & The sum of the four amounts\\
Average & B7 & The average of the four amounts\LotEnd
\end{LotTable}

\section{Interpreting the Results}
\LotExample{examples/expenses5}
The total is the sum of the amounts in B2..B5; the average is the total
divided by the number of expenses.  The dollar sign and the two decimal
places come from the format, not from what you typed.  An amount of a
thousand dollars or more needs a wider column: if a cell shows a row of
asterisks, widen the column with \slashcmd{/Worksheet Column Set-Width}.

\section{Changing the Template}
To change an amount, type the new number over the old one: the total and
the average follow it.

To add an expense:
\begin{Steps123}
\item Move the total and the average down one row: select
  \slashcmd{/Move}, type \type{a6..b7} and press \kname{ENTER}, type
  \type{a7} and press \kname{ENTER}.
\item Type the new expense in row 6: \type{Paper} in A6 and \type{35} in
  B6.
\item Give the new amount the Currency format: select
  \slashcmd{/Range Format Currency}, type \type{2} and press \kname{ENTER},
  type \type{b6} and press \kname{ENTER}.  The formats of the moved cells
  moved with them; the new cell has none of its own.
\item Move the cell pointer to B7, press \fkey{Edit}{F2}, press
  \kname{BACKSPACE} twice, type \type{6)} and press \kname{ENTER}.  Do the
  same in B8.  The moved formulas still added B2..B5; now they include
  B6.
\end{Steps123}
\FlairShot[fig:sa5b]{sa5b}{The report with a fifth expense}
